\hypertarget{ux5d0ux5e8ux5dbux5d9ux5d8ux5e7ux5d8ux5d5ux5e8ux5d4}{%
\section{ארכיטקטורה}\label{ux5d0ux5e8ux5dbux5d9ux5d8ux5e7ux5d8ux5d5ux5e8ux5d4}}

\hypertarget{ux5e7ux5e9ux5d1-ux5e0ux5e7ux5d5ux5d3ux5ea-ux5deux5dbux5e4ux5dcux5d4-ux5deux5d3ux5d5ux5e8ux5d2ux5ea}{%
\subsection{קשב נקודת-מכפלה
מדורגת}\label{ux5e7ux5e9ux5d1-ux5e0ux5e7ux5d5ux5d3ux5ea-ux5deux5dbux5e4ux5dcux5d4-ux5deux5d3ux5d5ux5e8ux5d2ux5ea}}

ארכיטקטורת ה-Transformer, שהוצגה לראשונה בעבודתם של ואסמאני ועמיתיו,
מבוססת על עיקרון יסודי של מנגנון קשב המאפשר לכל יסוד ברצף לשקלל את
השפעתם של יתר היסודות בעת בניית ייצוגו הפנימי. מנגנון זה, המכונה קשב
נקודת-מכפלה מדורגת (Scaled Dot-Product Attention), מהווה את אבן הפינה של
הארכיטקטורה כולה, והבנתו המדויקת חיונית לתפיסת כלל רבדי המודל.

הפעולה הבסיסית של מנגנון הקשב מגדירה שלושה ווקטורים עבור כל אסימון ברצף:
וקטור השאילתה (Query), וקטור המפתח (Key), ווקטור הערך (Value). שלושת
הווקטורים נגזרים מהייצוג הכניסה באמצעות מטריצות הטלה לינאריות נפרדות,
המסומנות בהתאמה כ-W\_Q, W\_K ו-W\_V. פירוק זה מאפשר למודל להבדיל בין
``מה אני מחפש'' (השאילתה), ``מה אני מציע'' (המפתח) ו''מה אני מעביר''
(הערך) --- הפרדה קונספטואלית המעניקה למנגנון את גמישותו.

פונקציית הקשב מוגדרת באופן הבא:

Attention(Q, K, V) = softmax( Q K\^{}T / sqrt(d\_k) ) V

כאשר d\_k מציין את ממדיות וקטורי המפתח והשאילתה. הגורם 1/sqrt(d\_k) מדרג
את מכפלות הנקודה כדי למנוע ריכוז הגרדיאנטים באזורים רוויים של פונקציית
ה-softmax, בעיה הנוצרת בממדים גבוהים שבהם מכפלות הנקודה נוטות לקחת ערכים
גדולים מאוד. פעולת ה-softmax מנרמלת את ציוני הקשב על פני כל האסימונים
ברצף, כך שסכומם שווה לאחד, ובכך הופכת אותם לפילוג הסתברות פשוט. התוצאה
הסופית היא שקלול משוקלל של וקטורי הערך לפי עוצמות הקשב, המשקפות את מידת
הרלוונטיות של כל אסימון עבור האסימון השואל.

ראוי לציין את קיומו של מנגנון מיסוך (Masking) המאפשר לחסום אסימונים
עתידיים בשכבות הדקודר, ולהסתיר אסימוני ריפוד (Padding) שנוספו לצורך
אחדות אורך בתוך אצווה. הדבר מושג באמצעות הוספת ערכים שליליים גדולים מאוד
(-∞) לאלמנטים הרצויים בטרם פעולת ה-softmax, ולמעשה מאפס את משקלם בפלט.

\hypertarget{ux5e7ux5e9ux5d1-ux5e8ux5d1-ux5e8ux5d0ux5e9ux5d9}{%
\subsection{קשב
רב-ראשי}\label{ux5e7ux5e9ux5d1-ux5e8ux5d1-ux5e8ux5d0ux5e9ux5d9}}

גם אם מנגנון קשב בודד מסוגל ללכוד תלויות לטווח ארוך בין אסימונים, הוא
מוגבל בכך שהוא מסנתז את כל המידע הרלוונטי לייצוג יחיד בכל צעד. מנגנון
הקשב הרב-ראשי (Multi-Head Attention) מרחיב גישה זו באמצעות הפעלה מקבילה
של מספר ראשי קשב עצמאיים, כל אחד עם מטריצות הטלה ייחודיות לו.

בפורמליזציה המתמטית, קשב רב-ראשי מוגדר כ:

MultiHead(Q, K, V) = Concat(head\_1, \ldots, head\_h) W\^{}O

כאשר כל ראש מוגדר על-ידי:

head\_i = Attention(Q W\_i\^{}Q, K W\_i\^{}K, V W\_i\^{}V)

מטריצות ה-W\_i\^{}Q, W\_i\^{}K, W\_i\^{}V הן מטריצות הטלה הניתנות
ללמידה, ייחודיות לכל ראש i, ומטריצת W\^{}O משמשת לקונקטנציה ולהטלה חוזרת
לממד המודל המלא d\_model. במימוש הנפוץ, ממדיות כל ראש נקבעת כ-d\_k =
d\_model / h, כך שמספר הפרמטרים הכולל דומה לקשב חד-ראשי ברוחב מלא.

החלוקה לראשים מרובים מאפשרת למודל ללמוד לכוון את תשומת לבו לאספקטים
שונים של הייצוג בו-זמנית: ראש אחד עשוי לשים דגש על תלויות תחביריות, אחר
על ייחוס ספרי (Coreference), ושלישי על יחסים סמנטיים רחבים יותר. ממצאים
אמפיריים מראים כי ראשים שונים של הקשב אכן מתמחים בסוגים שונים של תלויות,
ותופעה זו תועדה בהרחבה בניתוח ה-BERT. יכולת זו להסתכל בו-זמנית על
``תת-מרחבי ייצוג'' מרובים היא אחד היתרונות המרכזיים של Transformer על
פני ארכיטקטורות רקורנטיות.

\hypertarget{ux5e9ux5dbux5d1ux5d5ux5ea-ux5e7ux5d3ux5d9ux5deux5d4-ux5e0ux5d5ux5e8ux5deux5dcux5d9ux5d6ux5e6ux5d9ux5d4-ux5d5ux5d7ux5d9ux5d1ux5d5ux5e8ux5d9ux5dd-ux5e9ux5d9ux5d5ux5e8ux5d9ux5d9ux5dd}{%
\subsection{שכבות קדימה, נורמליזציה וחיבורים
שיוריים}\label{ux5e9ux5dbux5d1ux5d5ux5ea-ux5e7ux5d3ux5d9ux5deux5d4-ux5e0ux5d5ux5e8ux5deux5dcux5d9ux5d6ux5e6ux5d9ux5d4-ux5d5ux5d7ux5d9ux5d1ux5d5ux5e8ux5d9ux5dd-ux5e9ux5d9ux5d5ux5e8ux5d9ux5d9ux5dd}}

לכל שכבת Transformer שני תת-רכיבים עיקריים: שכבת הקשב הרב-ראשי ושכבת
קדימה מיקומית (Position-wise Feed-Forward Network). שכבת הקדימה מופעלת
על כל אסימון בנפרד ובאופן זהה, ומוגדרת כ:

FFN(x) = max(0, x W\_1 + b\_1) W\_2 + b\_2

כלומר, שתי טרנספורמציות לינאריות עם פונקציית הפעלה ReLU ביניהן. ממדיות
הרשת הפנימית d\_ff נקבעת לרוב לכ-4 כפול ממדיות המודל (d\_ff = 4 *
d\_model), מה שמאפשר לשכבה לבצע חישובים לא-לינאריים עשירים על כל ייצוג
אסימון בנפרד לאחר שהקשב אסף מידע מהרצף כולו.

לצד שני הרכיבים הללו, שני מנגנוני ייצוב אימון חיוניים שזורים
בארכיטקטורה: חיבורים שיוריים (Residual Connections) ונורמליזציית שכבה
(Layer Normalization). הגישה הראשונה, שמקורה ברשתות ResNet לעיבוד תמונה,
מוסיפה את קלט כל תת-שכבה לפלטה, כך שהרשת לומדת בפועל רק את ה''שארית''
(Residual) ביחס להתמרה הזהות:

x\_out = LayerNorm(x + Sublayer(x))

מנגנון זה מקל על זרימת הגרדיאנטים לאורך עומק הרשת ומאפשר אימון של ערימות
עמוקות מאוד. נורמליזציית השכבה (Layer Normalization) מנרמלת את הייצוגים
בתוך כל אסימון על פני ממדי המאפיין, שונה מנורמליזציית אצווה (Batch
Normalization) הפועלת על פני דוגמאות שונות. בחירה זו מתאימה במיוחד
לרצפים באורכים משתנים ולשימוש עם גדלי אצווה קטנים, ותורמת לייצוב ניכר של
תהליך האימון.

העומק המאפשר לצבור מספר שכבות כאלו --- במקור 6 שכבות בקודר ו-6 בדקודר
--- מאפשר לרשת לבנות ייצוגים הירארכיים עשירים, שבהם שכבות נמוכות לוכדות
תכונות לוקאליות ושכבות גבוהות מייצגות יחסים גלובאליים וסמנטיים מורכבים.

\hypertarget{ux5e7ux5d9ux5d3ux5d5ux5d3-ux5deux5d9ux5e7ux5d5ux5deux5d9}{%
\subsection{קידוד
מיקומי}\label{ux5e7ux5d9ux5d3ux5d5ux5d3-ux5deux5d9ux5e7ux5d5ux5deux5d9}}

בשונה ממודלים רקורנטיים המעבדים רצפים בסדר כרונולוגי ובכך שומרים מטבעם
על מידע סדרי, מנגנון הקשב הוא פעולה חסרת-סדר ביסודה --- תוצאתו אינה
משתנה אם האסימונים בקלט יסודרו מחדש. על מנת לאפשר למודל לנצל את מיקומם
של האסימונים ברצף, מוסיפים לוקטורי ההטמעה (Embedding) ייצוג מיקומי
מפורש.

בנוסחת הקידוד המיקומי המסינוסואידלי המקורית:

PE(pos, 2i) = sin(pos / 10000\^{}(2i / d\_model)) PE(pos, 2i+1) =
cos(pos / 10000\^{}(2i / d\_model))

כאשר pos הוא מיקום האסימון ברצף ו-i הוא מדד הממד. בחירה זו אינה
שרירותית: התכונות המתמטיות של פונקציות הסינוס והקוסינוס מבטיחות כי ניתן
לבטא את ה-PE של מיקום pos + k כשילוב לינארי של ה-PE של מיקום pos, ובכך
המודל יכול ללמוד להתייחס להיסטים יחסיים בנוסף למיקומים מוחלטים.

גישה חלופית ומאוחרת יותר היא שימוש בקידודים מיקומיים הניתנים ללמידה
(Learned Positional Encodings), שבה ייצוגי המיקום הם פרמטרים חופשיים
המתעדכנים במהלך האימון ממש כמו שאר הפרמטרים. גישה זו אומצה על-ידי BERT
וכן על-ידי ViT לתמונות. יתרונה הוא בגמישות --- המודל יכול לאמץ קידוד
מיקומי שמתאים בדיוק לנתונים --- אך חסרונה הוא בחוסר יכולת הכללה למיקומים
ארוכים יותר מאלו שנראו באימון.

גרסאות מודרניות של קידוד מיקומי, כגון קידוד מיקומי יחסי רוטציוני (RoPE,
Rotary Position Embedding) ו-ALiBi, שיפרו עוד יותר את יכולת ההכללה
לאורכי רצף שלא נצפו בזמן האימון, ופרטיהן נדונים בהרחבה בפרקים המאוחרים
יותר של ספר זה.
